\gddchapter{Analysis}

\section{Thematic Analysis Framework}
Six phase model to move from theoretical raw audio to theoretical insights:
This is just a scaffold that will be filled out with actual data from this session:


\begin{enumerate}
    \item Familiarisation
    
    Listen to the recorings while reading chapter 5 to notice intiial patterns

    \item Initial Coding
    
    Generate labels for specific UX moments such as "Command frustration", "modal confusion" etc.

    \item Searching for Themes
    
    Group codes into broader categories like "Neutral language friction" or "Immersion breaking"

    \item Reviewing Themes
    
    Check if those found themes accurately represent the difference between two input modalities, 
    and if they fit the main research question.

    \item Defining Themes
    
    Refine the essence of what the data tells you about your research themes

    \item Writing Up
    
    Integrate transcript extracts in the finial presentation that show how the player 
    interacted with the experience, including the themes defined above as anchoring points.

\end{enumerate}

\section{Sticky moments analysis}

Analyse the moments of interst noted earlier in Chapter 3. Remember not to treat the system misunderstanding the player as a faliure.
Instead, note the general patterns that came up and try to generate insights from them.

\section{Language fluency consideration}

Include the level of English language that the participant used and their percieved degree of fluency.
Are they native speakers? What about their accent? Make sure to include that too.